\documentclass[12pt]{article}
\usepackage[T1]{fontenc}
\usepackage[utf8]{inputenc}
\usepackage{lmodern}
\usepackage{textcomp}
\usepackage{enumitem}
\usepackage{geometry}
\geometry{margin=1in}
\begin{document}
\begin{center}
Answer Key - Biology - Undergraduate Level Assessment 11 \
\small (Answers are ordered exactly as in the question paper.)
\end{center}

\section*{Multiple Choice}
\begin{enumerate}[label=\arabic*.]
\item \textbf{Answer:} D
\\ \textit{Options:} A) Flash flood; B) Tsunami; C) Sudden reduction in the animal food resource; D) Destruction of the ozone layer; E) Drought
\\ \textit{Note:} The destruction of the ozone layer is a biotic factor that affects the growth rate of a population because it leads to increased UV radiation, which can harm living organisms and disrupt ecosystems, ultimately impacting population dynamics.
\item \textbf{Answer:} C
\\ \textit{Options:} A) It is an organic compound made of proteins.; B) It is a catalyst that alters the rate of a reaction.; C) It is operative over a wide pH range.; D) The rate of catalysis is affected by the concentration of substrate.
\\ \textit{Note:} Trypsin is an enzyme that functions optimally at a specific pH range (around 7.5 to 8.5) and is not operative over a wide pH range, making the statement that it is operative over a wide pH range untrue.
\item \textbf{Answer:} D
\\ \textit{Options:} A) Cause euphoria and hallucinations; B) Improve memory and cognitive functions; C) Increase physical strength and endurance; D) Cause tremors, spasms, or even death; E) Trigger temporary paralysis without pain
\\ \textit{Note:} The breakdown of cholinesterase by nerve gases leads to an accumulation of acetylcholine at neuromuscular junctions, resulting in continuous stimulation of muscles, which causes tremors, spasms, and can ultimately lead to respiratory failure and death.
\item \textbf{Answer:} A
\\ \textit{Options:} A) Enzymatic reactions with oxygen; B) Lactose crystallization; C) Loss of refrigeration power over time; D) mesophiles; E) Natural separation of milk components without bacterial influence
\\ \textit{Note:} Milk spoils in the refrigerator primarily due to enzymatic reactions facilitated by residual enzymes and microbial activity, which can still occur at lower temperatures, leading to the breakdown of proteins and fats when exposed to oxygen.
\item \textbf{Answer:} D
\\ \textit{Options:} A) the potential mates experience geographic isolation; B) the potential mates experience behavioral isolation; C) the potential mates have different courtship rituals; D) the potential mates have similar breeding seasons
\\ \textit{Note:} The correct answer is that similar breeding seasons among potential mates can actually facilitate interspecific breeding rather than prevent it, as it allows for the possibility of mating during the same time frame.
\end{enumerate}

\section*{True / False}
\begin{enumerate}[label=\arabic*.]
\setcounter{enumi}{5}
\item \textbf{Answer:} False
\\ \textit{Options:} A) True; B) False
\\ \textit{Note:} The statement is false because the correct probability of drawing an ace first and a face card second, without replacement, is calculated as (4/52) * (12/51), which simplifies to 48/2652, or 1/55, not 4/221.
\item \textbf{Answer:} True
\\ \textit{Options:} A) True; B) False
\\ \textit{Note:} Polytene chromosomes in Drosophila silvarentis are formed by the replication of a single chromosome without cell division, resulting in multiple identical chromatid strands that can be observed under a light microscope.
\item \textbf{Answer:} True
\\ \textit{Options:} A) True; B) False
\\ \textit{Note:} Animal cells contain centrioles within centrosomes, which play a crucial role in organizing microtubules and facilitating the separation of chromosomes during cell division.
\item \textbf{Answer:} False
\\ \textit{Options:} A) True; B) False
\\ \textit{Note:} Deserts have low rates of leaching and nutrient cycling due to their limited vegetation and extremely low rainfall, which restricts soil moisture and biological activity essential for these processes.
\item \textbf{Answer:} False
\\ \textit{Options:} A) True; B) False
\\ \textit{Note:} Sexual dimorphism is primarily a result of natural selection and evolutionary pressures related to reproduction and survival, rather than artificial selection imposed by humans.
\end{enumerate}

\section*{Long Form Response}
\begin{enumerate}[label=\arabic*.]
\setcounter{enumi}{10}
\item \textbf{Answer:} In the mating of a black male cat and a tortoiseshell female, the expected phenotypic ratios of the offspring can be understood through the principles of sex-linked inheritance. Tortoiseshell coloration in cats is linked to the X chromosome, with the black coat color typically being dominant. The offspring can be categorized as follows: (1/4) tortoiseshell females, who inherit one X chromosome with the black allele and one with the orange allele; (1/4) black females, who inherit two X chromosomes with the black allele; (1/4) yellow males, who inherit the X chromosome with the orange allele and a Y chromosome; and (1/4) black males, who inherit the black allele from their mother and a Y chromosome from their father. This ratio reflects the genetic contributions from both parents and the sex-linked nature of the coat color genes.
\\ \textit{Note:} correct because it accurately applies the principles of sex-linked inheritance to predict the phenotypic ratios of the offspring, distinguishing between the genotypes of male and female cats based on their X chromosome contributions.
\item \textbf{Answer:} Arthropods exhibit key characteristics that have significantly contributed to their evolutionary success, particularly their development of a chitinous exoskeleton, well- developed organ systems, and complex organization of nervous tissues. The chitinous exoskeleton provides structural support and protection, allowing arthropods to thrive in diverse environments while also facilitating mobility through jointed appendages. Additionally, their advanced organ systems, including specialized respiratory and circulatory systems, enable efficient gas exchange and nutrient distribution, which are crucial for survival in various habitats. The complexity of their nervous system allows for enhanced sensory perception and rapid responses to environmental stimuli, further aiding in their adaptation to changing conditions. Collectively, these characteristics underscore the evolutionary advantages that have positioned arthropods as one of the most successful groups of organisms on Earth.
\\ \textit{Note:} The answer correctly identifies the chitinous exoskeleton, advanced organ systems, and jointed appendages as key characteristics that enhance the adaptability and survival of arthropods in various environments, demonstrating their evolutionary success.
\end{enumerate}

\vfill
\noindent\textit{End of Answer Key}
\end{document}